\documentclass[11pt,a4paper]{article}

% PAQUETES
\usepackage[T1]{fontenc}%
\usepackage[utf8]{inputenc}%

\usepackage[english]{babel}
\usepackage{amsmath}
\usepackage{amsthm}
\usepackage{amsfonts}
\usepackage[%left=1.54cm,right=1.54cm,top=1.54cm,bottom=1.54cm
    margin=1in,
]{geometry}
\usepackage{xfrac}  
\usepackage{tikz-cd}
\usepackage{enumerate}
\usepackage{amsfonts}
\usepackage{amssymb}
\usepackage{tcolorbox}
\usepackage{rotating}
\usepackage{mathpazo}
% \usepackage{charter}
\usetikzlibrary{babel}
\usepackage{listings}
\usepackage{amssymb}
\usepackage{extarrows}
\usepackage{makeidx}
\usepackage{graphicx}
\usepackage{multirow}
\usepackage{tikz-cd}
\usepackage{tasks}
\usepackage{xcolor}

% Christian
\usepackage{enumitem,etoolbox,titlesec}


%OPERADORES

\DeclareMathOperator{\dom}{dom}
\DeclareMathOperator{\cod}{cod}
\DeclareMathOperator{\id}{id}

\newcommand{\red}[1]{\textcolor{red}{#1}}
\newcommand{\C}{\mathbb{C}}
\newcommand{\im}{\text{Im}}
\newcommand{\R}{\mathbb{R}}
\newcommand{\N}{\mathbb{N}}
\newcommand{\Z}{\mathbb{Z}}
\newcommand{\D}{\mathbb{D}}
\newcommand{\B}{\mbox{Ob}}
\newcommand{\M}{\mbox{Mo}}
\newcommand{\del}{\Delta}
\newcommand{\odel}[1]{\left[#1\right]}
\newcommand{\Hom}[1]{\text{Hom}(#1)}
\newcommand{\adel}[1]{\left\lbrace #1 \right\rbrace}
%\renewcommand{\theequation}{\thesection.\arabic{equation}}
\newcommand{\funcion}[5]{%
{\setlength{\arraycolsep}{2pt}
\begin{array}{r@{}ccl}
#1\colon \hspace{0pt}& #2 & \longrightarrow & #3\\
& #4 & \longmapsto & #5
\end{array}}}

\newcommand{\func}[3]{#1\colon  #2  \to  #3}

\newcommand\restr[2]{{% we make the whole thing an ordinary symbol
  \left.\kern-\nulldelimiterspace % automatically resize the bar with \right
  #1 % the function
  \vphantom{\big|} % pretend it's a little taller at normal size
  \right|_{#2} % this is the delimiter
  }}

%ENTORNOS

% \theoremstyle{theorem}
\newtheorem{teo}{Theorem}[section]
\newtheorem{prop}[teo]{Proposition}
\newtheorem{lem}[teo]{Lemma}
\newtheorem{cor}[teo]{Corollary}

\theoremstyle{definition}
\newtheorem{defi}[teo]{Definition}
\newtheorem{rem}{Remark}[teo]
\newtheorem{exa}{Example}
\newtheorem{eje}{Exercise}
\newtheorem{que}{Question}

\newenvironment{sol}
  {\begin{proof}[\textit{Solution}]}
  {\end{proof}}

\titlelabel{\thetitle.\quad}

\def\contador{}
\graphicspath{{./figures/}}
\newcommand{\qand}{\quad\text{and}\quad}
\usepackage{microtype,parskip}
\setlength{\parindent}{15pt}
\linespread{1.15}
\usepackage{hyperref}
\hypersetup{
    colorlinks=true,  
    allcolors=blue,
    pdfproducer={Christian Chávez},
}

\makeatletter
\@ifclassloaded{exam}{
    \footer{}{\thepage}{}
    \renewcommand{\thequestion}{\bfseries\arabic{question}}
    \renewcommand{\solutiontitle}{\noindent\textit{Solution.}\enspace}
    \unframedsolutions
}{}
\makeatother


\newenvironment{theproof}
{
    \renewcommand{\solutiontitle}{}
    \begin{solution}
    \vspace*{-\baselineskip}
    \begin{proof}
}
{
    \end{proof}
    \end{solution}
    \renewcommand{\solutiontitle}{\noindent\textit{Solution.} }
}

\begin{document}

\def\contador{Lesson 5}
\noindent
\begin{minipage}[c]{0.33\textwidth}
    \includegraphics[scale=0.37]{sello_yachay.png}
\end{minipage}
\begin{minipage}[c]{0.37\textwidth}
    % \centering
    \textbf{\large School of Mathematical and\\ Computational Sciences}\par
    Abstract Algebra
\end{minipage}
~ 
\begin{minipage}[c]{3mm}
    \raggedleft
    \rule[1.5mm]{0.3mm}{15mm}
\end{minipage}
~ 
\begin{minipage}[c]{0.24\textwidth}
    \raggedleft
    Prof. Pablo Rosero\\
    \& Christian Chávez\\
    \contador{}
\end{minipage}

\vspace{1mm}
\noindent\hrulefill

\vspace{3mm}

In this lecture, we meet two concepts that help us study groups and their properties.
The objective is not only  to learn this concepts, but mainly to  use them to study the principal objects of group theory.

\section{Homomorphisms and isomorphisms}

When studying an algebraic structure, we are not only interested in the objects that possess this structure, but also in the morphisms between objects of the same type that preserve such structure. 
This is true through out all branches of mathematics.

In the case of group theory, the objects are groups and  the structure preserving maps are called group-homomorphisms.

Let \((G,\cdot)\) and \((H, *) \) be groups.

\begin{defi}
    A group-homomorphism from \(G\) to \(H\) is a map \(\varphi \colon G\to H\) such that 
    \[\varphi(a\cdot b) = \varphi(a) * \varphi(b)\]
    for all \(a,b\in G\).
\end{defi}

\begin{rem}
    From now and on, we omit the adjective ``group'' in ``group-homomorphism'', since we will be working only on the realm of group theory \((\mathbf{Grp})\).
\end{rem}

Note that \(\varphi\) transforms a product \(a\cdot b\) (using the operation of \(G\))  into the   product \(\varphi(a)*\varphi(b)\) (using the operation of \(H\)).
This is why we say \(\varphi\) preserves the structure: it takes a product in \(G\) and maps it to a product in \(H\).
In this case (of groups) there are no more structure to be preserved.

However, the definition above just guarantees that the structure is preserved   in one direction only: from \(G\) to \(H\).
If, in addition, it is possible to preserve 
the structure the other way around, we obtain a more interesting type of map: an \textit{isomorphism}.

\begin{defi}
    An \textbf{isomorphism} is a bijective homomorphism.
\end{defi}

Thus, an isomorphism \(\varphi\colon G\to H\) is just a homomorphism that has   inverse \(\varphi^{-1}\colon H\to G\).
In this case, \(G\) and \(H\) are said to be of the same isomorphism type, or \textit{isomorphic} for short. We write \(G\cong H\).
You may ask \textit{isn't  it necessary that  the inverse   be a homomorphism also (so that the structure is preserved the other way around)?}
And it turns out that \textit{no}, because it follows straight from the definition. (Why?)

\begin{eje}
    Prove \(\varphi^{-1}\) is a homomorphism if \(\varphi\) is an isomorphism.
\end{eje}

\begin{eje}
    Prove \(\cong\) is an equivalence relation (over which set?)
\end{eje}


\begin{lem}\label{lem:facts.pf.group.isomorphism}
    If \(\varphi\colon G\to H\) is an isomorphism, then 
    \begin{enumerate}[label=(\roman*)]
        \item \(|G| = |H|\)
        \item \(G\) is Abelian if and only if \(H\) is Abelian
        \item \(\varphi\) preserves the order of elements, that is, \(|x| = |\varphi(x)|\).
    \end{enumerate}
\end{lem}

\begin{exa}
\begin{enumerate}[label=(\roman*)]
    \item Let us prove \(\left( \R, +  \right)\cong \left( \R^+, \cdot \right)\).
    Define \(\psi\colon \R\to \R^+:x\mapsto \exp(x)\). 
    Notice \[\psi(x + y) = \exp(x+y) = \exp(x)\cdot\exp(y) = \psi(x)\cdot\psi(y).\]
    Since \(\psi\) is injective and surjective (facts known from elementary calculus), \(\psi\) is an isomorphism.
    (Are there any other maps that show \(\left( \R, +  \right)\cong \left( \R^+, \cdot \right)\)?)
    \item If \(X\) and \(Y\) have the same cardinality, then \(S_X\cong S_Y\).
    Indeed, if \(\omega \colon X\to Y\) is a bijection (recall the definition of cardinality), and
    given any \(\alpha\in S_X\), 
    the map \(\omega \circ \alpha\circ \omega^{-1}\colon Y\to Y \) is also a bijection.
    Thus, define \(\varphi\colon S_X\to S_Y\) by \(\varphi(\alpha) = \omega \circ \alpha\circ \omega^{-1}\) for every \(\alpha\in S_X\).
    Diagramatically, we have
    \[\begin{tikzcd}
X \arrow[r, "\alpha", two heads, hook] & X \arrow[r, "\omega", two heads, hook] & Y \arrow[r, "\omega\circ \alpha\circ \omega^{-1}", two heads, hook] &[3em] Y
\end{tikzcd}.\]
    Since \[\varphi(\alpha\circ \beta) = \omega \circ (\alpha\circ\beta)\circ \omega^{-1}= (\omega \circ \alpha\circ \omega^{-1})\circ(\omega\circ\beta\circ \omega^{-1}) = \varphi(\alpha)\circ \varphi(\beta),\]
    \(\varphi\) is a homomorphism.
    The reader must easily verify that \(\varphi\) is an isomorphism.

    \item The groups \(S_3\) and \(\Z/6\Z\) are not isomorphic as one is Abelian and the other is not. This follows from Theorem \ref{lem:facts.pf.group.isomorphism}.
    
    \item \(\left( \R\setminus \left\{ 0 \right\}, \cdot \right) \) and \(\left( \R,+ \right)\) are not isomorphic.
    The former  has an element of  order \(2\) and the latter does not have any element of order \(2\).
\end{enumerate}
\end{exa}

\begin{eje}
\begin{enumerate}[label=(\roman*)]
    \item Provide an example of a group with only one element.
    \item Prove \(S_3\) is the unique nonabelian group of order \(6\) up to isomorphism.
    \item Prove \(A\times B\cong B\times A\) if \(A\) and \(B\) are any groups.
\end{enumerate}
\end{eje}



\section{Subgroups}




We now come to study of groups that live inside larger groups.

Let \((G,\cdot_G)\) be a group.

\begin{defi}
    A group \((H,\cdot_H)\) is a \textbf{subgroup} of    \((G,\cdot_G)\) if \(H\subseteq G\) and \(\cdot_H = \cdot_G\).
    
\end{defi}

In other words, a subgroup of \(G\) is a subset of \(G\) that it together with the operation of \(G\) is itself a group.
If \(H\) is a subgroup of \(G\) we write \(H\leq G\).
As an exercise, prove \(\leq \) is a partial order.

\begin{exa}
\begin{enumerate}[label=(\roman*)]
    \item Under addition, \(\Z\leq  \mathbb{Q} \leq  \R \).
    \item There are always two noninteresting subgroups of any group \(G\), namely \(\left\{ 1_G \right\}\) and \(G\) itself.
    We call the former the \textit{trivial subgroup} of \(G\).
    \item If $G=D_{2 n}$ is the dihedral group of order $2 n$, let $H$ be $\left\{1, r, r^2, \ldots, r^{n-1}\right\}$, the set of all rotations in $G$. Since the product of two rotations is again a rotation and the inverse of a rotation is also a rotation it follows that $H$ is a subgroup of $D_{2 n}$ of order $n$.
\end{enumerate}
\end{exa}


The following result says that the subgroups of \(G\) are precisely the nonempty subsets of \(G\) that are closed under multiplication and inverses.

\begin{lem}[Subgroup criterion]\label{lem:subgroup.criterion}
    A set \(H\) is a subgroup of \(G\) if and only if 
    \begin{enumerate}[label=(\roman*)]
        \item \(\varnothing\neq  H \subseteq G \), and 
        \item \(ab^{-1}\in H\) for all \(a,b\in H\).
    \end{enumerate}
\end{lem}


\begin{proof}
\begin{enumerate}[label=(\roman*)]
    \item[(\(\Rightarrow\))] This follows from the definition.
    \item[(\(\Leftarrow\))] Suppose \((i)\) and \((ii)\) hold.
    Asociativity of the multiplication of \(G\) holds for any subset of \(G\), so in particular it holds on \(H\).
    Because \(H\) is nonempty, we can choose \(x\in H\) and using the second condition we see \(1_G = x x^{-1} \in H\).
    Further, if \(a\) is any element of \(H\),  \((ii)\) implies \(a^{-1} = 1_G a^{-1}\in H\).
    Thus \(H\) is closed under inverses.
    Finally, note \(H\) is closed under products since \(a b = a (b^{-1})^{-1}\in H\) for any \(a,b\in H\).

\end{enumerate}
\end{proof}

\begin{eje}
    Prove that if \(H\) is a nonempty finite subgroup of \(G\) that is closed under multiplication, then \(H\) is a subgroup of \(G\).
\end{eje}

\subsection{Some important subgroups}



\begin{defi}
    Let $G$ be a group and  $A \subseteq G$.
    \begin{enumerate}[label=(\roman*)]
        \item \(C_G(A) = \{g \in G \mid g a g^{-1} = a \text{ for all } a \in A\}\) is the \textbf{centralizer} of $A$.
        \item \(Z(G) = \{g \in G \mid g a = a g \text{ for all } a \in G\}\) is the \textbf{center} of $G$.
        \item \(N_G(A) = \{g \in G \mid g A g^{-1} = A\}
\) is the \textbf{normalizer} of $A$, where $g A g^{-1} = \{g x g^{-1} \mid x \in A\}$.
 
    \end{enumerate}
\end{defi}

In words, the {centralizer} of a subset \( A \) in a group \( G \) is the set of elements in \( G \) that commute with all elements of \( A \). 
The center of \(G\) is the centralizer of \(G\).
The {normalizer} is the set of elements that leave \( A \) fixed under conjugation. Both the centralizer and the normalizer are subgroups of \( G \).

\begin{prop}\hfill\null
    \begin{enumerate}[label=(\roman*)]
        \item \(C_G(A) , N_G(A), Z(G) \leq G\)
        \item \(C_G(G) = G\)
        \item \(C_G(A)\leq N_G(A)\)
    \end{enumerate}
\end{prop}

\begin{proof}
\begin{enumerate}[label=(\roman*)]
    \item Let us prove \(N_G(A) \leq G\).
    First note \(N_G(A)\neq \varnothing\) as \(1\in N_G(A)\).
    Now let \(x,y\in N_G(A)\).
    We want to prove that \(xy^{-1}\in N_G(A)\), i.e., 
    \[(xy^{-1})A(xy^{-1})^{-1} = A\quad\text{or}\quad xy^{-1}Ayx^{-1} = A.\]
    Note \(gAg^{-1} = A\) if and only if \(g^{-1}Ag=A\).
    Thus we have \(y^{-1}Ay = A\).
    Now the equality above reduces to 
    \(xAx^{-1}=A\). This is  true because \(x\in N_G(A)\).
    Hence we reach our goal.
    By the subgroup criterion (Lemma \ref{lem:subgroup.criterion}), we conclude \(N_G(A)\leq G\).


    \item The rest of the proof is trivial and is left to the reader.
\end{enumerate}
\end{proof}

\begin{eje}[Classwork]
    Assume the following to be true.
\begin{enumerate}[label=(\roman*)]
    \item If $G$ is a group and $N \trianglelefteq G$, then $|H| \mid |G|$.
    \item If $H \leq G$, then $H \leq N_G(H)$.
    \item If $H \trianglelefteq G$, then $H = C_G(H)$ if and only if $H$ is abelian.
\end{enumerate}
Prove that if $G = S_3$ and $A = \{(1), (1\ 2\ 3), (1\ 3\ 2)\}$, then $C_G(A) = A$ and $N_G(A) = G$.
\end{eje}
\end{document}